% Chapter 9: settings.
\chapter{Settings}\label{ch:settings}
\index{settings}

The machine keeps seven numbers in its own memory. They survive power-off.
You change them only from a PC with the \textbf{winder} program
(Chapter~\ref{ch:pc}); there is no settings menu on the panel. Most owners
never need to touch them: calibration (Chapter~\ref{ch:calibration}) sets the
important ones for you.

The machine counts length in \emph{encoder counts}. From the factory, one
inch of paper is about 191 counts, and one count is about 0.005 inch. The
program converts to inches and yards for you.

\begin{center}
\small
\renewcommand{\arraystretch}{1.45}
\begin{tabularx}{\textwidth}{@{}>{\raggedright\arraybackslash}p{0.27\textwidth}>{\raggedright\arraybackslash}p{0.22\textwidth}>{\raggedright\arraybackslash}X@{}}
\toprule
\textbf{Setting (name in the program / short name)} & \textbf{Factory value} & \textbf{What it does, allowed range} \\
\midrule
\ui{Measuring wheel size}\par \setting{wheel\_um} &
  4.0015 in across\par (101638 micrometres) &
  Wheel diameter. Turns counts into inches. Range 50000 to 300000 micrometres
  (about 2 to 11.8 in). \\
\ui{Encoder counts per wheel turn}\par \setting{cpr} &
  2400 &
  Counts per full turn of the wheel. Must match the encoder (600 lines x 4).
  Range 100 to 100000. \\
\ui{Short wind (30 yd button)}\par \setting{short} &
  30.0 yd\par (206187 counts) &
  Target length for \btn{1}. Range 1000 to 4000000 counts. \\
\ui{Long wind (90 yd button)}\par \setting{long} &
  90.0 yd\par (618562 counts) &
  Target length for \btn{2}. Same range. \\
\ui{Coast after the motor stops}\par \setting{coast} &
  72 in (2 yd)\par (13746 counts) &
  How far the roll glides after the motor is switched off. The motor is cut
  this far before the goal. Must be smaller than both targets. Set by
  calibration. Saved settings keep their coast after a firmware update:
  use \ui{Reset to factory settings} and Save to get the 2 yd default. \\
\emph{(not in the program)}\par \setting{margin} &
  1.0 in\par (191 counts) &
  Extra length added to every roll for the guarantee. Range 0 to 2000 counts
  (about 0 to 10.5 in). \\
\emph{(not in the program)}\par \setting{bump} &
  1 (on) &
  1 = trim mode with the minimum-amount guarantee. 0 = legacy mode (no trim,
  no guarantee). \\
\bottomrule
\end{tabularx}
\end{center}

One more rule: the longer target plus the margin may not exceed 4000000
counts (about 580 yards). The machine refuses settings that break any rule.

\section{Changing a setting}
\begin{enumerate}
  \item Start the \textbf{winder} program and open the \ui{Settings} tab.
  \item Type the new value. A banner \ui{Unsaved changes} appears.
  \item Click \ui{Save to the machine}. Nothing changes on the machine until
        you do. \ui{Undo all} throws your edits away.
\end{enumerate}
\ui{Reset to factory settings} puts every value back to the table above.
Check the screen afterwards: if the \ui{Unsaved changes} banner shows, click
\ui{Save to the machine} to store them.

\begin{warning}[Change the wheel and counts only with care]
A wrong \ui{Measuring wheel size} or \ui{Encoder counts per wheel turn}
makes every roll the wrong length. After changing either, run a calibration
with the length check (Chapter~\ref{ch:calibration}).
\end{warning}

\section{Margin and legacy mode}\label{sec:set-bump}
\index{legacy mode}
The \setting{margin} and \setting{bump} settings are not shown in the program
and calibration never changes them. They can only be changed through the
serial command line by a technician
(see \path{docs/hardware/serial-protocol.md} in the source repository).
Leave \setting{bump} at 1. With \setting{bump} at 0 the machine behaves like
the older firmware: it cuts the motor once and accepts wherever the roll
stops, with no trim, no \shows{UNDER TARGET} check and no guarantee.

\section{When settings are lost}\label{sec:settings-lost}
\index{settings!lost}
The settings are stored with a check code. If the stored settings are
damaged, or were saved by firmware older than fw=6, the machine ignores them
and uses the factory values. After updating to fw=\fwver{} from an older
version you must calibrate (or re-enter your numbers) and save once. The
program's \ui{Connection check} reports this as \ui{Saved settings are valid}
failing.
